The following conditions are equivalent:
\[
(i)\ S\in\mathfrak A_n,
\qquad
(ii)\ \kappa(S)<c_n,
\qquad
(iii)\ \gamma(S)<1.
\]
When these conditions hold, \(\rho^*(S)=\sqrt{\gamma(S)}\); otherwise \(\rho^*(S)=1\).  Let \(p=\dim(S)\), let \(\ell_i=(P_S)_{ii}\), and put \(\beta=p/(n-1)\).  For \(t\in[\beta,1]\), define
\[
d_i(t)=\frac{n-1}{n}t-\ell_i
\]
and
\[
U(\ell)=\min\left\{t\in[\beta,1]:d_i(t)\geq0\ \text{for every }i,\quad
2\max_i\sqrt{d_i(t)}\leq\sum_{i=1}^n\sqrt{d_i(t)}\right\}.
\]
Then the critical squared quality is determined exactly by the score leverage profile:
\[
\gamma(S)
=\min_{\substack{K\in\mathbb S^n:\ K\mathbf1_n=0,\\
\operatorname{diag}(K)=\operatorname{diag}(P_S)}}
\lambda_{\max}(K\mid H_n)
=U(\ell).
\]
In particular, the sharp bounds are
\[
\beta\leq\frac{n}{n-1}\max_i\ell_i\leq\gamma(S)\leq1,
\]
and
\[
\gamma(S)=\beta
\quad\Longleftrightarrow\quad
\operatorname{diag}(P_S)=\frac{p}{n}(1,\ldots,1)^\top.
\]
For the three rank-one four-unit score spaces in the endpoint proposition, the formula gives respectively
\[
\gamma(S_4)=\frac13,
\qquad
\gamma(S_{\mathrm{mid}})=\frac89,
\qquad
\gamma(S_{\mathrm{out}})=1.
\]
Consequently,
\[
\{\rho^*(S)\}^2\geq\frac{p}{n-1},
\qquad
V^*(\rho;S)<V_{\mathrm{CRD}}(\rho;S)
\Longrightarrow
\rho^2>\gamma(S)\geq\frac{n}{n-1}\max_i\ell_i.
\]
Given a basis matrix for \(S\), forming \(P_S\) and solving the displayed one-dimensional problem computes \(\gamma(S)\), decides membership in \(\mathfrak A_n\), and computes \(\rho^*(S)\) from the score matrix alone, without enumerating balanced assignments.  For a specified \(\rho\), enumeration of \(\mathcal Z_n\) and the finite semidefinite program in the exact oracle saddle theorem compute \(v^*(\rho;S)\), \(\kappa(S)\), and an \(\varepsilon\)-optimal attaining assignment mixture.